\documentclass[a4paper,14pt]{extarticle}
\usepackage{styles}

\begin{document}

\litSubsubsection{3.1.1 Experiment Conduction and Benchmarks}

This section includes the experimental validation of stated hypothesis - GPU acceleration of genetic algorithms execution. The validation strategy uses deliberately simplified ``toy problems'' (well-established benchmarks in the computational optimization literature) - XOR perceptron training and the Traveling Salesman Problem (TSP).

\textbf{Toy problems isolation advantages.} A critical advantage of toy problems for GPU acceleration validation is their ability to separate hardware acceleration effects from algorithmic improvements. Deliberately unoptimized genetic algorithm implementations are used - they ensure that measured performance differences reflect pure GPU acceleration capabilities. This makes the experiment applicable across variations of fitness functions as black-box optimization objects.

\textbf{Toy Problem Selection.} The use of simplified benchmark problems for performance evaluation is a common methodological approach in computer science and optimization research. Such toy problems serve to isolate specific algorithmic or implementation characteristics to controllable, reproducible test cases. This approach allows to measure pure acceleration effects without taking into account variables introduced by problem-specific optimizations or domain.

The XOR problem has been extensively used as a benchmark for evaluating neural network training algorithms~\cite{googlebooks_xor_perceptron,schaffer2001genetic, samir2017xor} and genetic algorithm implementations. It represents a simple non-linear problem and is considered a benchmark for testing neural network capabilities in solving more complex problems. It is linearly inseparable and this task makes it suitable for genetic algorithm evaluation: the problem is being sufficiently challenging to require evolutionary search, but simple enough to provide consistent, reproducible results to measure performance. Additional experimental advantage: it is deterministic and the genetic algorithm is executing a fixed number of generations regardless of the test results. This prevents early termination that would introduce variance in timing measurements and ensures each experimental run performs equivalent computational work.

Similarly, the Traveling Salesman Problem serves as one of the most widely recognized NP-complete problems and benchmarks in combinatorial optimization tasks. It is considered one of ``high-dimensional problem instances with a known global optimal solution, that are very well suited as benchmark problems for metaheuristics as for example GAs''~\cite{affenzeller2009genetic}. TSP is used in the literature to evaluate parallel computing performance because its fitness function (path cost calculation) is computationally intensive, but also highly parallelizable. With $O(n^2)$ complexity of the distance matrix and each candidate solution as a dot product, this problem is a representation of a typical CUDA-accelerated matrix-matrix multiplication task~\cite{kirk2016programming}.

\subsubsection*{Optimization Methods}

\textbf{Shared Memory for Read-Only Data.} Training generational data (XOR inputs/outputs, TSP distance matrices) is cached in GPU shared memory to exploit the latency reduction compared to global memory. As documented in NVIDIA's architecture guides~\cite{nvidia2011cuda} and validated through microbenchmarking studies~\cite{wong2010demystifying}, memory hierarchy latencies differ by an order of magnitude:

\begin{itemize}
\item Global Memory Access: 400-800 cycles
\item Shared Memory Access: 20-40 cycles
\end{itemize}

Implementation requires cooperative loading at block initialization, followed by \texttt{\_\_syncthreads()} barrier synchronization to ensure data visibility before shared access throughout kernel execution. Wong et al.~\cite{wong2012cuda} are one of the researchers to demonstrate that shared memory utilization is important for achieving optimal GPU performance in compute- and memory-bound algorithms.     

\textbf{Population Residency (Persistent Device Buffers).} Traditional genetic algorithm implementations transfer the entire population between CPU and GPU at each generation, incurring $O(\text{generations})$ transfer overhead. Research on GPU-accelerated evolutionary algorithms~\cite{li2018accelerating,luong2013gpu} consistently identifies memory transfer reduction as critical to performance.

In this paper an optimized approach was tested to maintain population residency on the GPU throughout evolution:

\textbf{Traditional:}
\begin{verbatim}
For each generation:
  CPU → GPU (population)
  GPU kernel (fitness)
  GPU → CPU (fitnesses)
  CPU: selection, crossover, mutation
\end{verbatim}

\textbf{Optimized:}
\begin{verbatim}
Initial: CPU → GPU (population, training data, RNG state)
For each generation:
  GPU kernel (fitness)
  GPU → CPU (fitnesses)
  CPU: selection via argsort
  CPU → GPU (sorted indices)
  GPU kernel (crossover + mutation via GPU RNG)
Final: GPU → CPU (best solution)
\end{verbatim}

This optimization keeps the population and training data GPU-resident, eliminating the per-generation population transfer. Per-generation transfers are reduced to fitness scores and sorted indices (small vectors of size $N$), rather than the full population matrix.

\textbf{Pinned (Page-Locked) Host Memory.} Standard host memory allocations are pageable
and require the GPU driver to copy data into temporary page-locked buffers before DMA
transfer~\cite{harris2012optimize,mao2020pinned}. This introduces additional latency and
reduces effective bandwidth. As NVIDIA documents, ``page-locked or pinned memory transfers
attain the highest bandwidth between the host and the device. On PCIe x16 Gen3 cards, for
example, pinned memory can attain roughly 12~GB/s transfer rates''~\cite{nvidia2024cuda}.
Pinned memory enables direct DMA access:
\begin{verbatim}
# Standard allocation (hidden staging copy)
arr = np.zeros(size)
cuda.to_device(arr)  # pageable → pinned staging → DMA

# Pinned allocation (direct DMA)
arr = cuda.pinned_array(size)
cuda.to_device(arr) # direct DMA transfer
\end{verbatim}

Figure~\ref{fig:host_device_transfer} illustrates the two transfer paths. The pageable path requires an intermediate CPU-mediated copy into a driver-managed staging buffer before DMA can begin, while the pinned path enables direct DMA access over the PCIe bus.

\begin{figure}[H]
\centering
\includegraphics[width=0.75\textwidth]{diagrams/sec2/host_device_transfer.png}
\caption{Host--device memory transfer paths: pageable vs.\ pinned allocation. Adapted from~\cite{nvidia2024cuda,sanders2010cuda}.}
\label{fig:host_device_transfer}
\end{figure}

Pinned memory reduces per-transfer latency by eliminating the intermediate staging copy, increasing effective PCIe bandwidth.
As NVIDIA's optimization guide states, ``batching many small transfers into one larger
transfer performs significantly better than making each transfer separately, even if
doing so requires packing non-contiguous regions of memory into a contiguous buffer and
then unpacking after the transfer''~\cite{nvidia2024cuda}.

\textbf{GPU-Side Random Number Generation.} CPU-generated random numbers require
transferring data every generation and create synchronization stalls that idle the GPU
pipeline. Each transfer adds fixed latency overhead regardless of size, while CPU-GPU
sync prevents effective latency hiding. This is a structural consequence of CUDA's
execution model: as NVIDIA specifies, ``memory copy, memory set functions, and kernel
calls that use the default stream begin only after all preceding calls on the device
(in any stream) have completed, and no operation on the device (in any stream) commences
until they are finished''~\cite{nvidia2024cuda}.

Using GPU-native random number generation (\texttt{numba.cuda.random}, implementing Xoroshiro128p~\cite{xoroshiro}):

\textbf{CPU approach:} generate $\rightarrow$ transfer $\rightarrow$ GPU waits\\
\textbf{GPU approach:} generate in-kernel $\rightarrow$ immediate use

Xoroshiro128p~\cite{xoroshiro} is a high-quality, fast PRNG that produces deterministic, reproducible sequences when initialized with a fixed seed per thread. Since each CUDA thread maintains its own independent RNG state in private registers, global memory contention is avoided entirely - ensuring both statistical independence across parallel threads and thread-safety without race conditions, even under massive parallelism.

This reduces transfer latency, removes synchronization barriers, and prevents pipeline bubbles (GPU idle periods while awaiting data), while maintaining reproducibility for scientific validation.

\textbf{Fundamental Principle: Latency vs. Bandwidth.} For small transfers (kilobyte-scale), latency dominates total transfer time. "PCIe (Peripheral Component Interconnect Express) is the
defacto interconnect for high-performance intra-host communications." ~\cite{hou2024pcie}. It is the high-speed serial bus interface used for CPU--GPU data transfers, where each ``lane'' provides a dedicated point-to-point link and multiple lanes operate in parallel. PCIe transfer characteristics follow:
\[
\text{Transfer time} = \text{latency} + \left(\frac{\text{size}}{\text{bandwidth}}\right)
\]

\subsubsection*{Development Tooling}

\begin{center}
\textbf{Table 1: GPU Computing and Data Processing Tools}\\[0.5em]
\begin{tabular}{@{}ll@{}}
\toprule
\textbf{Tool} & \textbf{Purpose} \\
\midrule
Numba CUDA & Python-to-CUDA JIT compilation \\
NumPy & Host-side array operations and data preparation \\
CuPy & GPU array operations (NumPy-compatible API) \\
\bottomrule
\end{tabular}
\end{center}

\begin{center}
\textbf{Table 2: Hardware Configuration}\\[0.5em]
\begin{tabular}{@{}ll@{}}
\toprule
\textbf{Component} & \textbf{Specification} \\
\midrule
GPU & NVIDIA GeForce RTX 2050 (4GB VRAM), 16 SMs, CC 8.6 \\
CPU & 12-core Intel processor, 16GB RAM \\
\bottomrule
\end{tabular}
\end{center}

\begin{center}
\textbf{Table 3: Statistical Analysis Tools}\\[0.5em]
\small
\begin{tabular}{@{}ll@{}}
\toprule
\textbf{Tool} & \textbf{Purpose} \\
\midrule
SciPy & Paired t-test, Wilcoxon signed-rank, Fisher's exact \\
NumPy & Descriptive statistics, bootstrap resampling for CIs \\
Matplotlib & Performance results visualization \\
\bottomrule
\end{tabular}
\end{center}

\begin{center}
\textbf{Table 4: Profiling and Performance Analysis Tools}\\[0.5em]
\small
\begin{tabular}{@{}ll@{}}
\toprule
\textbf{Tool} & \textbf{Purpose} \\
\midrule
Numba \texttt{@jit(nopython=True)} & Kernel compilation with execution timing \\
\texttt{numba.cuda.profiler} & GPU timing regions and kernel statistics \\
CuPy Profiler & Array operation timing and memory usage \\
Nsight Compute (NCU) & Kernel metrics (occupancy, throughput) \\
Nsight Systems (NSYS) & System-wide timeline (CPU/GPU/transfers) \\
\bottomrule
\end{tabular}
\end{center}


\subsubsection*{Statistical Analysis}

\textbf{Sample Size and Power Analysis.} Scientific validity in performance evaluation requires sufficient statistical analysis and relevant sample sizes to detect meaningful effects with adequate statistical power. Power analysis (as established in Cohen's framework~\cite{cohen1988power}) was used as a primary method to determine the minimum number of experimental runs needed to achieve 80\% probability of detecting true effects at $\alpha=0.05$ significance:

\begin{center}
\begin{tabular}{@{}lr@{}}
\toprule
\textbf{Effect Size (Cohen's $d$)} & \textbf{Required $N$ ($\alpha=0.05$, power=0.80)} \\
\midrule
Small (0.2) & 393 \\
Medium (0.5) & 63 \\
Large (0.8) & 25 \\
\bottomrule
\end{tabular}
\end{center}
    
For the following experiments of this course paper, $N=30$ independent runs per configuration provides sufficient power to detect large effects while at the same time maintaining computational feasibility for making claims derived from the results.

\textbf{Statistical Tests.} To perform direct comparison of GPU versus CPU performance, paired experimental design with identical random seeds was used:

\begin{center}
\small
\begin{tabular}{@{}lll@{}}
\toprule
\textbf{Comparison} & \textbf{Test} & \textbf{Rationale} \\
\midrule
Execution time & Paired t-test & Paired samples, normal timing data \\
Fitness values & Wilcoxon signed-rank & Non-parametric, discrete fitness scores \\
\bottomrule
\end{tabular}
\end{center}

\textbf{Reported Metrics.} Performance analysis includes: (1) speedup factor as the ratio of CPU to GPU wall-clock time, (2) 95\% confidence intervals computed via bootstrap resampling for ratio estimators~\cite{efron1994bootstrap}, (3) p-values at $\alpha=0.05$ significance, and (4) time decomposition to separate computation from memory transfer overhead.


\end{document}